\chapter{Die Idee: Zahlen, die wissen, wie viel Platz sie haben}
\label{ch:idee}

% Ebene: Oberfläche. Keine abgesetzten Formeln, nur Zahlen im Fließtext.
% Entstehungsreihenfolge beachten: umgedrehte Basen -> Endlichkeit ->
% führende Nullen (Korrektur vom 13.07.2026).

\section{Eine Ordnung, die niemand hinterfragt}

Jedes Stellenwertsystem beruht auf einer stillen Vereinbarung. Im
Dezimalsystem zählt die hinterste Ziffer Einer, die davor Zehner, dann
Hunderter -- jede Stelle wiegt das Zehnfache ihrer Nachbarin. Im
Binärsystem dasselbe mit Zweierschritten. Und es gibt ein weniger
bekanntes, aber altehrwürdiges System, in dem die Basen nicht konstant
sind, sondern wachsen: das Fakultätssystem. Dort darf die hinterste
Ziffer nur 0 oder 1 sein, die nächste 0 bis 2, die übernächste 0 bis 3,
und so weiter. Je weiter links eine Stelle steht, desto mehr darf sie
zählen \emph{und} desto mehr ist sie wert.

In allen diesen Systemen zeigt dieselbe Richtung nach oben: Die Stelle
mit dem meisten Spielraum trägt das meiste Gewicht. Das wirkt so
selbstverständlich, dass man es nie ausspricht.

Die Horimetrik beginnt mit der Frage, was passiert, wenn man genau das
umdreht. Was, wenn ausgerechnet die \emph{engste} Stelle -- die, die
nur 0 oder 1 kennt -- das größte Gewicht bekommt? Wenn die Stellen
nach rechts hin immer großzügiger werden, ihr Einfluss aber immer
kleiner?

\section{Die erzwungene Endlichkeit}

Die erste Konsequenz kommt sofort, und sie ist keine Designentscheidung,
sondern ein Zwang. Ein gewöhnliches Zahlensystem wächst nach links ins
Unendliche: Wird eine Zahl zu groß, stellt man eine weitere Stelle
voran, mit noch größerem Gewicht -- Platz ist immer. Im umgedrehten
System geht das nicht. Links steht bereits die engste Stelle, die mit
der Basis zwei, und unterhalb von zwei gibt es nichts mehr. Das System
ist nach links \emph{versiegelt}.

Wer mit einer festen Stellenzahl arbeitet, bekommt deshalb einen
endlichen Vorrat an Zahlen. Mit einer Stelle sind es zwei, mit zweien
sechs, mit dreien vierundzwanzig, mit vieren hundertzwanzig -- jede
zusätzliche Stelle multipliziert den Vorrat, und die Größen, die dabei
entstehen, sind genau die Fakultäten. Ist der Vorrat aufgebraucht,
hilft kein Übertrag: Man braucht ein \emph{neues System} mit mehr
Stellen, in dem jede alte Ziffer plötzlich ein anderes Gewicht trägt.

Diesen endlichen Lebensraum einer Zahl nennen wir ihren
\emph{Horizont}. Was zunächst wie ein Defekt des umgedrehten Systems
aussieht -- man kann nicht einfach weiterzählen --, ist in Wahrheit
sein interessantester Zug: Jede Zahl trägt hier nicht nur einen Wert,
sondern auch die Information, in welchem Raum dieser Wert lebt. Die
Gesamtheit aller Horizonte bleibt dabei unendlich; nur jeder einzelne
ist endlich. Unendlichkeit gibt es weiterhin -- aber als Treppe, nicht
als Gerade.

\section{Der zweite Fund: die Nullen}

Die führenden Nullen kamen erst später, und sie kamen als Überraschung.
In jedem gewohnten System ist eine vorangestellte Null reine Kosmetik:
012 ist 12. Im umgedrehten System stellt man eine Null voran -- und
verändert damit die Zahl.

Ein Beispiel mit kleinen Zahlen. Im Horizont mit zwei Stellen bedeutet
die Ziffernfolge 12 den Wert fünf. Stellt man eine Null davor, steht
dort 012 -- aber jetzt in einem Horizont mit drei Stellen, in dem alle
Gewichte anders verteilt sind. Der Wert ist sechs. Noch eine Null:
0012, Wert sieben. Die Ziffern 1 und 2 haben sich nicht bewegt, und
doch zählt die Zahl bei jeder Verlängerung eins weiter.

Das ist der Moment, in dem aus einer Kuriosität ein
Untersuchungsgegenstand wird. Denn wenn eine führende Null den Wert
verändert, dann gibt es plötzlich \emph{zwei verschiedene Antworten}
auf die harmloseste aller Fragen: Was heißt es, einer Zahl mehr Platz
zu geben?

\begin{itemize}
  \item Man kann ihren \emph{Wert} festhalten -- dann müssen sich beim
        Umzug in den größeren Horizont ihre Ziffern umorganisieren.
  \item Man kann ihre \emph{Gestalt} festhalten -- die Ziffern bleiben
        stehen, eine Null kommt davor -- dann verschiebt sich ihr Wert.
\end{itemize}

Im Dezimalsystem sind das dieselbe Operation. Hier sind es zwei, und
sie führen in verschiedene Welten. Fast alles, was dieses Buch
untersucht, lebt in der Lücke zwischen diesen beiden Begriffen: die
beiden zunächst sichtbaren Identitäten einer Hori-Zahl (eine dritte
wird sich später zeigen), die Arten aufzuräumen, die Rechenwelten
und die Spannungen zwischen ihnen.

\section{Was hier eigentlich erfunden wurde -- und was nicht}

Ehrlichkeit gehört von Anfang an dazu: Die Bausteine dieses Systems
sind sämtlich bekannte Mathematik. Zahlensysteme mit gemischten Basen
gibt es seit Jahrhunderten, das Fakultätssystem ist gut untersucht,
und dass seine Räume Fakultätsgröße haben, weiß die Kombinatorik
lange. Neu -- jedenfalls haben wir es in dieser Form nirgends gefunden
-- ist die Kombination: die umgedrehte Leserichtung in einem endlichen
Raum, der Horizont als Bestandteil der Zahl, die bedeutungstragenden
Nullen und die beiden konkurrierenden Arten zu wachsen. Ob diese
Kombination eine eigenständige Theorie trägt, war die Wette, mit der
dieses Projekt begann. Die folgenden Kapitel legen die Karten offen.

\begin{oberflaeche}
Falls von diesem Kapitel nur ein Satz bleiben soll, dann dieser: Die
Horimetrik entstand nicht aus dem Wunsch, Zahlen anders zu schreiben,
sondern aus der Frage, was eine Ordnung erzwingt, wenn man sie
umdreht. Die Antwort war: Endlichkeit -- und aus der Endlichkeit
folgte alles Weitere von selbst.
\end{oberflaeche}
